\documentclass[12pt]{article}
\usepackage{setspace}
\usepackage{amsmath,amssymb,euscript,graphicx}
\usepackage{xurl}
\usepackage{graphicx}
\usepackage{subfigure}
\usepackage{wrapfig}
\usepackage[compact]{titlesec}
\usepackage{fancybox,color}
\usepackage[footnotesize]{caption}
\usepackage{cite}
\usepackage{enumitem}
\usepackage[normalem]{ulem}
\usepackage[top=0.9375in, right=0.9375in, left=0.9375in, bottom=0.9375in, centering]{geometry}
\usepackage{bm}

\usepackage{amsthm}


% -- Loading the code block package:
\usepackage{listings}
% -- Basic formatting
\usepackage[utf8]{inputenc}
\usepackage[english]{babel}
\usepackage{times}
\setlength{\parindent}{8pt}
\usepackage{indentfirst}
% -- Defining colors:
\usepackage[dvipsnames]{xcolor}
\definecolor{codegreen}{rgb}{0,0.6,0}
\definecolor{codegray}{rgb}{0.5,0.5,0.5}
\definecolor{codepurple}{rgb}{0.58,0,0.82}
\definecolor{backcolour}{rgb}{0.95,0.95,0.92}
% Definig a custom style:
\lstdefinestyle{mystyle}{
    backgroundcolor=\color{backcolour},   
    commentstyle=\color{codepurple},
    keywordstyle=\color{NavyBlue},
    numberstyle=\tiny\color{codegray},
    stringstyle=\color{codepurple},
    basicstyle=\ttfamily\footnotesize\bfseries,
    breakatwhitespace=false,         
    breaklines=true,                 
    captionpos=t,                    
    keepspaces=true,                 
    numbers=left,                    
    numbersep=5pt,                  
    showspaces=false,                
    showstringspaces=false,
    showtabs=false,                  
    tabsize=2
}
% -- Setting up the custom style:
\lstset{style=mystyle}


% IEEE uses Times Roman font, so we'll default to Times.
% These three commands make up the entire times.sty package.
\renewcommand{\sfdefault}{phv}
\renewcommand{\rmdefault}{ptm}
\renewcommand{\ttdefault}{pcr}

\newcommand{\squeeze}{\vspace{-0.1in}}
\newcommand{\changed}[1]{\textcolor{blue}{#1}} %generates text in blue

%%%%%%%%%%%%%%%%%%%%%%%%%%%%%%%%%%%%%%%%%%%%%%%%%%%%%%%%%%%%%%%%%%%%%%%

\begin{document}

\begin{center}
  {\Large {\bf Lab 1: Turtle Graphics and Terminal Commands}}
  
\end{center}

\vspace{5pt}
\hrule
\vspace{5pt}

\section{Summary}
In this lab, we will introduce basic terminal functionality and Turtle Graphics.

\section{Introduction to the Terminal}
We will go over what a terminal is, how to use it, and why you should care.

\subsection{Open the Terminal}
Open the VS code text editor.

Open the terminal by either clicking \texttt{Ctrl + `} or clicking \texttt{Terminal} in the menu and selecting \texttt{New Terminal}. Note: In Windows this is at the top of the window and in Mac this is in the menu bar. If you do not see it, then try selecting the 3 dots in the menu bar and then selecting Terminal $\rightarrow$ New Terminal. 

\subsection{What is the Terminal?}
The terminal is a text-based interface to everything on your computer that allows you to manipulate files, execute computer programs, and open documents.
Unlike your graphical interface like Finder (Mac) or Explorer (Windows) which allow you to point and click different icons and drag things around the screen using the mouse, a terminal only allows input from the keyboard.
Anything you can do using your point and click interface you can do using a terminal, however, the reverse is not true.

Note that historically Windows uses a command prompt terminal interface and Mac OS / Linux use a bash terminal interface.
There are differences in the commands used in these systems and for consistency we have set up everyone's VS Code to use a bash style interface. 


\subsubsection{Why you should care?}
\begin{itemize}
\item Terminal is a long standing tool in computing that has well established tools that are not changing or going away. 
\item There are easy ways to do tasks with lots of files, \textit{e.g.}, renaming 1000 files with a single line in the terminal.
\item Finally, the code we will write this semester will run in the terminal.
\end{itemize}


\subsection{Directories and Paths}
\subsubsection{How does a terminal work?}

Just like in your Finder or Explorer interface, when you are in a specific folder this corresponds to a \textbf{directory} in the terminal.
The current directory in a terminal is called the \textbf{working directory}.

A terminal will initially open to the home directory of your computer account.
For example on a Mac,

\vspace{5mm}
\textit{/Users/$<$username$>$}
\vspace{5mm}
where username is the computer account name.

Similarly, on a Windows it will look something like,

\vspace{5mm}
\textit{C:\textbackslash Users\textbackslash $<$username$>$}
\vspace{5mm}
where username is the computer account name.

\textbf{Note to Windows users}: instructions use the forward slash convention moving forward. On windows in the Terminal you can use backwards and forwards slashes interchangeably. This is not true on Mac. 

The directory represents the place in computer memory where the information is stored.
The complete description of where a directory or file on a computer is located is called a \textbf{path}.
Note that paths are \textbf{case-sensitive}. 
We will use paths and filenames to change directories, list contents of directories, and run Python programs. 

\subsubsection{If I cannot click in the terminal how do I change directories?}

As a text-based interface, we need key stroke commands to move directories the way you would click through folders in Finder or Explorer.

\begin{enumerate}
\item \textit{pwd}: Stands for ``print working directory''. It tells you where you currently are on your computer. For example when you open the terminal, usually it is in the home folder. Test this out, type \textit{pwd} in the terminal and see what comes out. 
\item \textit{ls}: Stands for ``list'' the files/folders in the current directory. 
\item \textit{mkdir $<$ new folder name here $>$}: Stands for ``make directory''. When you type this command followed by a space and then the name of a new folder you want to create in your current working directory. 
  \item \textit{cd $<$ type in a directory here $>$}: Stands for ``change directory''. When you type this command followed by a space, followed by the directory path you want to go to the Terminal directory will change. 
\end{enumerate}

\subsubsection{Do you always need to know the exact path?}
Can you instead navigate the Terminal like you do a point-and-click interface?

\begin{center}
  \begin{tabular}{| l | l | }
    \hline
    \textbf{Command} & \textbf{Explanation} \\ \hline
    cd        & Set the current directory to your home directory  \\ \hline
    cd \textasciitilde{}     & \shortstack[l]{Set the current directory to your home directory. \\ (The \textasciitilde{} is shorthand for your home directory)} \\  \hline
    cd ..     & \shortstack[l]{Move up to the parent directory (or one folder up) from the \\ current working directory} \\ \hline
    cd ../../ & \shortstack[l]{Move up to the parent of the parent directory (or two folders up)\\ from the current working directory} \\ \hline
    cd blah   & \shortstack[l]{Move to the sub-directory blah (Note this must be a directory\\ in your current working directory)} \\ \hline
    cd /Users/$<$username $>$/Desktop & Move to the Desktop directory given by the absolute path \\ \hline
    cd - & Move to your last working directory \\ \hline
\end{tabular}
\end{center}


\subsubsection{Additional Useful Terminal Commands}
Beyond changing directories and looking at files in those directories there are many other helpful terminal commands. 

For example
\begin{itemize}
\item \textit{mv $<$from$>$ $<$to$>$}: move a file from one location to another location. This includes just renaming it in the current directory. The mv command will remove the file in the old location after copying it to the new location.
\item \textit{cp $<$from$>$ $<$to$>$}: copy a file to a new location or name
\item \textit{rm $<$filename$>$}: remove the file (will not remove a directory). \textbf{WARNING:} This is not the same as dragging a file to the trash in Finder or Explorer. Once the file is removed there is no recovering the file.
\item \textit{rm -r $<$directory$>$}: remove the directory and all the files and subdirectories in it.
\item \textit{less $<$filename $>$}: scroll through a file
\item \textit{cat $<$filename$>$}: send the file to standard output in the terminal
\item \textit{echo ``a string''}: write the text within the string to standard output in the terminal 
  \item \textit{touch $<$filename$>$}:updates the modification date on a file or creates an empty file if the named file does not exist. This can be useful in various situations, like when you are learning to create, rename, and delete files using the terminal
\end{itemize}

\subsection{Terminal Tips and Tricks}
There are some helpful features and best practices to use when learning the terminal
\begin{itemize}
\item Use \textit{pwd} and \textit{ls} often as you navigate through directories to see what directory you are in and what files are in the directory.
\item You can hit return anywhere in a terminal line and it will execute the entire line no matter where the cursor is located. Note you cannot navigate in the line using a mouse, you have to use the arrow keys or other builtin shorthands or Command actions (See Table \ref{table:command_action})
  \item Tab-Completion: you can type part of a name of a file or program, and then press the tab key to complete the filename as far as possible while the choice is unique. If two or more files are similarly named the terminal will output the options of files to continue with. 
  \item Dragging a folder to change directories: While you are learning about changing directories it is possible to make folders in your Finder or Explorer windows and drag these folders to the terminal. Type cd + a space in the terminal and then drag a folder to the terminal. The absolute path of that folder will appear in the terminal. Hit enter to navigate to this location. 
    \item The terminal has builtin manual pages also called man pages. If you forget what a particular terminal page might do you can type \textit{man pwd} and a manual page will come up detailing the complete functionality of the command (which has far more capability than we cover in this lab). To exit a man page type 'q'. 
\end{itemize}

\begin{table}[h]
  \centering
  \begin{tabular}{| l | l | }
    \hline
    \textbf{Command} & \textbf{Explanation} \\ \hline
    cntl-a & move the cursor to the start of the line \\ \hline
    cntl-b & move the cursor back one position \\ \hline
    cntl-c & kill the current process running in the terminal \\ \hline
    cntl-d & delete the character under the cursor \\ \hline
    cntl-e & move the cursor to the end of the current line \\ \hline
    cntl-f & move the cursor forward one position \\ \hline
    cntl-k & delete everything on the current line to the right of the cursor \\ \hline
    cntl-l & clear the terminal screen \\ \hline
    cntl-n & move forward in the terminal history \\ \hline
    cntl-p & move backward in terminal history \\ \hline
    cntl-z & \shortstack[l]{freeze the current process running in a termial. \\ You can put it in the background by typing \textit{bg} \\or put it in the foreground by typing \textit{fg}}\\ \hline
  \end{tabular}
  \caption{Command Actions}
  \label{table:command_action}
\end{table}

\subsection{Create a lab1 folder inside of your ES1098 folder}
To showcase your understanding of this section, you will now create a lab1 folder and move into that directory. 
\begin{itemize}
  \item open a terminal in VS Code. 
  \item Navigate to your ES1098 folder from lab0, use \texttt{ls}, \texttt{cd}, and \texttt{pwd} to help you find this directory
  \item Once you are in this directory you should see: \texttt{/Users/<computername>/ES1098} when you type \texttt{pwd}.
  \item Create a new folder by typing \texttt{mkdir lab1}
  \item Move into that folder by typing \texttt{cd lab1}
  \item Confirm your current directory by typing \texttt{pwd}.
\end{itemize}

\section{Turtle Graphics}
Python is a powerful interpreted language that can do many different things. Last fall we learned the same concepts you will this spring but did assignments that manipulated and created data. This spring we will focus on art and sound for our assignments.
To this end, we will use turtle graphics to draw images in our computers.
You can image that there is a turtle holding a paint brush on your screen and when given commands the turtle will move leaving paint behind as it goes.
The commands you will need for this lab include: 

\begin{itemize}
\item \texttt{forward(x)} - move forward \texttt{x} pixels
\item \texttt{backward(x)} - move backward \texttt{x} pixels
\item \texttt{left(a)} - turn left \texttt{a} degrees
\item \texttt{right(a)} - turn right \texttt{a} degrees
\item \texttt{up()} - pick up the pen (don't draw when the turtle moves)
\item \texttt{down()} - put down the pen down (draw when the turtle moves)
\end{itemize}

For more turtle details consult the documentation:

\url{https://docs.python.org/3/library/turtle.html}

\subsection{Make a new python program}
You cannot write code in Microsoft or Google Docs. Instead you write code in a text editor that removes fancy fonts and other formatting features.

\begin{itemize}
\item Open VS Code
  \item We will create a file called \texttt{hexagon.py}. 
  \item Create a new file, either using File $\rightarrow$ new file, For Mac: \texttt{Cmd + N}. For Windows: \texttt{Ctrl + N}. Or from the terminal type \texttt{code hexagon.py}.
  \item At the top of every file from now on you will include a header that is not read by the python interpreter but is useful for grading. Please include the following in your file:
\begin{lstlisting}[language=Python]
"""
    Author: <YOUR NAME>
    Date: <TODAY'S DATE>    
    Purpose: What your file does
"""\end{lstlisting}
\item The first line of code will import turtle graphics:
  \begin{lstlisting}[language=Python]
from turtle import * \end{lstlisting}
\item The last line will be
    \begin{lstlisting}[language=Python]
exitonclick() \end{lstlisting}
    This will keep the turtle window open until you click on the window that pops up
  \item Save your file \texttt{Cmd + S} for MAC or \texttt{Ctrl + S} for Windows. If prompted supply the name \texttt{hexagon.py}. The \texttt{.py} tells your computer that you are writing Python code. 
  \item Using the \texttt{right()} and \texttt{forward()} commands draw a hexagon. One line of code will go on each line between the import and exitonclick() commands.
  \item To run your program you will type in your terminal:

    \vspace{5mm}
    \textit{python hexagon.py}
    \vspace{5mm}

  \item If something does not work, before asking a question make sure to save your file and run it again in the terminal.
  \item You should not move on until you see a hexagon on the screen.
\end{itemize}

\subsubsection{Things didn't work, what do I do?}
In computer science, when things do not work as you intend there is most likely a bug in your code.
Some bugs are syntax bugs meaning that you have an issue with how your code is written and it cannot be interpreted by the python interpreter.
Some bugs are logic bugs meaning that your code will run because nothing is wrong with the syntax but it is not doing what you wanted it to do.

If your code from above is not working here are some things to try / consider:
\begin{itemize}
\item Did you get an error in your terminal? Error messages are very specific in the terminal and often can point you to the problem and a potential solution. Check that everything is spelled correctly, that all parenthesis are closed, and that there are no spaces or indentations before any command. (These are syntax errors)
\item Are you seeing lines on the screen but it is not the shape that you want? (This is a logic bug). Consider this, how many sides are in a hexagon? Remember from geometry that the sum of all exterior angles is 360 degrees. Using this information, how can you determine the number of forward and right commands you need? 
\end{itemize}

To help a bit more here is an example of an L shape.
\begin{lstlisting}[language=Python]
"""
  Author: Victoria Edwards
  Date: 03/18/2026
  Purpose: Example for lab 1
"""
from turtle import *

forward(100)
right(90)
forward(50)

exitonclick() \end{lstlisting}

\subsubsection{Why is \texttt{exitonclick()} the last command?}
Comment out \texttt{exitonclick()} using the pound key (\#) before the line of code.

\begin{itemize}
  \item What happens?
  \item Why do you need this line of code at the end of the file? 
\end{itemize}

\section{Submission}
For each project you will include a brief written report along with your code files.
We are going to go over the submission procedure now as a class. 

\subsection{Report}
Each report consists of:
\begin{enumerate}
\item Header at the top of the page with your name, date and class
\item A title that describes the content of the report, for example, ``Victoria Edwards' ES 1098 Lab 1: My Hexagon''. 
\item Abstract: A one paragraph summary of the project
\item Results: Images showing what you created alongside a description of what each image is and how you designed / created it with code.
\item Conclusion: A brief description (1-3 sentences) of what you learned.
\item AI Usage Statement: A statement on if you used AI on this project. Reflection on how you used AI for the project including the software used. If you did not use AI for the project you still need to include a usage statement that says you did not use the tool. Include reflection on how you might have used the tool. Include all prompts that you used and how the results were incorporated into the project. A good rule of thumb, is to only submit work that you can completely explain. If you use code you do not understand / cannot explain this is considered a violation of the academic integrity policy.
\item Acknowledgments: A list of websites that you used and people you worked with including TAs and professors.
\end{enumerate}

You will create a google doc attached to the assignment in google classroom and fill it in with this material for lab 1. 

\subsection{Code submission}
Each week you need to include all code that is necessary to create the results shown in your report.
Please do not delete code, instead comment code out as necessary.
If you have results in your report and there is no corresponding code I have no way of knowing where you got the result from.

In google classroom, upload your \texttt{hexagon.py} file.

Remember to submit your lab project. 


\section{Acknowledgments}
Section 2, 3, and 4 of these instructions are from Professor Oliver Layton and Professor Hannen Wolfe in the Colby College CS Department. 
Sources:
\begin{itemize}
\item \footnotesize{\url{https://cs.colby.edu/courses/S25/cs151/projects/00/Lab00Software.html}}
\item \url{https://cs.colby.edu/courses/S25/cs151/projects/01/Lab01First-Python.html}
\item \footnotesize{\url{https://stackoverflow.com/questions/42606837/how-do-i-use-bash-on-windows-from-the-visual-studio-code-integrated-terminal}}
\end{itemize}


\end{document}
